\par\noindent\textbf{(1)} 输入谱为Ã(Ω,0)∝exp(−T0²Ω²/2)，传播相位增加k1zΩ+k2zΩ²/2。定义随群移动的时间τ=t−k1z后，频域积分核只剩exp[−(T0²−ik2z)Ω²/2−iΩτ]，说明脉冲中心群延迟为k1z。
\par\noindent\textbf{(2)} 使用∫exp(−aΩ²/2−iΩτ)dΩ∝a\textasciicircum{}\{−1/2\}exp(−τ²/(2a))，取D=k2z/T0²，可得A(z,τ)=A0(1−iD)\textasciicircum{}\{−1/2\}exp[−τ²/(2T0²(1−iD))]，此式同时保留幅度归一与色散相位。
\par\noindent\textbf{(3)} 取模方得到|A|²=|A0|²(1+D²)\textasciicircum{}\{−1/2\}exp[−τ²/(T0²(1+D²))]，故强度1/e半宽为T=T0√(1+D²)。指数虚部相位为−Dτ²/[2T0²(1+D²)]，其时间导数正比τ，形成近线性频率啁啾。
\par\noindent\textbf{(4)} k2→0时D→0，只剩中心时移而无展宽。|D|≫1时T≈|k2|z/T0，脉冲越短谱越宽，在固定色散距离上可能展宽越显著；非线性自相位调制、吸收和高阶色散都会破坏上述线性高斯解。
